% !TEX program = context
% !TEX root = mechanism-chapter-3.tex
%
% Mary Somerville, Mechanism of the Heavens (1831).
% Book I, Chapter III: On the Equilibrium of a System of Bodies.
% Original printed pages 54--68; source PDF pages 130--144.
% Run here: context mechanism-chapter-3.tex
% Styling, paper color and texture controls are in mechanism-environment.tex.
% Original wording and notation are retained; comments identify source errors.
% Figures 34--40 are schematic MetaFun reconstructions.

\environment mechanism-environment
\setupinteraction[state=stop,
  title={Mechanism of the Heavens - Book I, Chapter III},author={Mary Somerville}]
\ifdefined\MechanismCombined\else\setupuserpagenumber[number=54]\fi
\HeavensChapterHeaders{III}{\ifodd\userpageno A SYSTEM OF BODIES.\else ON THE EQUILIBRIUM OF\fi}

% Keep the original equation numbers and chapter-specific reference names.
\unexpanded\def\BookFormula#1#2{%
  \setnumber[formula][\number\numexpr#1-1\relax]%
  \placeformula[eq:chapter3-#1]%
  \startformula #2\stopformula
}

\ifdefined\MechanismCombined\else\starttext\fi

\blank[12mm]
\midaligned{\tfa CHAPTER III.}
\blank[4mm]
\midaligned{\sc ON THE EQUILIBRIUM OF A SYSTEM OF BODIES.}
\blank[4mm]

% Book I, Chapter III, articles 114--126; printed pages 54--59.
% Source PDF pages 130--135. Figures are schematic MetaFun reconstructions.

\startuseMPgraphic{heavens:figure34}
  numeric u; u := .23mm;
  pair M,N,P,A;
  P := (0,0); M := (143u,0); A := (59u,0); N := (59u,10u);
  pickup pencircle scaled .35pt;
  draw P--N--M; draw P--M; draw N--A;
  label.lft(textext("\tfx\m{m'}"),P);
  label.rt(textext("\tfx\m{m}"),M);
  label.top(textext("\tfx\m{n}"),N);
  label.bot(textext("\tfx\m{a}"),A);
  label(textext("\tfx{\it fig.}\,34."),(70u,28u));
\stopuseMPgraphic

\startuseMPgraphic{heavens:figure35}
  numeric u; u := .23mm;
  pair S,M,P,Q,N,NN,NNN,A,B,C,AA,BB,CC,AAA,BBB,CCC;
  S := (0,0); M := (77u,53u); P := (102u,-43u); Q := (155u,0);
  N := (61u,27u); NN := (86u,-28u); NNN := (128u,0);
  A := whatever[S,M]; (N-A) dotprod (M-S) = 0;
  B := whatever[M,P]; (N-B) dotprod (P-M) = 0;
  C := whatever[M,Q]; (N-C) dotprod (Q-M) = 0;
  AA := whatever[S,P]; (NN-AA) dotprod (P-S) = 0;
  BB := whatever[P,M]; (NN-BB) dotprod (M-P) = 0;
  CC := whatever[P,Q]; (NN-CC) dotprod (Q-P) = 0;
  AAA := NNN;
  BBB := whatever[Q,P]; (NNN-BBB) dotprod (P-Q) = 0;
  CCC := whatever[Q,M]; (NNN-CCC) dotprod (M-Q) = 0;
  pickup pencircle scaled .35pt;
  draw S--M--Q--P--S; draw S--Q; draw M--P;
  draw M--N; draw N--A; draw N--B; draw N--C;
  draw P--NN; draw NN--AA; draw NN--BB; draw NN--CC;
  draw Q--NNN; draw NNN--BBB; draw NNN--CCC;
  label.lft(textext("\tfx S"),S);
  label.top(textext("\tfx\m{m}"),M);
  label.bot(textext("\tfx\m{m'}"),P+(0,-7u));
  label.rt(textext("\tfx\m{m''}"),Q);
  label.bot(textext("\tfx\m{n}"),N);
  label.ulft(textext("\tfx\m{n'}"),NN);
  label.urt(textext("\tfx\m{a}"),A);
  label.rt(textext("\tfx\m{b}"),B);
  label.rt(textext("\tfx\m{c}"),C);
  label.bot(textext("\tfx\m{a'}"),AA);
  label.urt(textext("\tfx\m{b'}"),BB);
  label.lrt(textext("\tfx\m{c'}"),CC+(7u,0));
  label.ulft(textext("\tfx\m{a''}"),AAA);
  label.lrt(textext("\tfx\m{b''}"),BBB);
  label.urt(textext("\tfx\m{c''}"),CCC);
  label(textext("\tfx{\it fig.}\,35."),(17u,50u));
\stopuseMPgraphic

% Original printed page 54.
\blank[medium]
\midaligned{\it Definitions and Axioms.}
\blank[small]

\Article{114}Any number of bodies which can in any way mutually affect each other's motion or rest, is a system of bodies.

\Article{115}Momentum is the product of the mass and the velocity of a body.

\Article{116}Force is proportional to velocity, and momentum is proportional to the product of the velocity and the mass; hence the only difference between the equilibrium of a particle and that of a solid body is, that a particle is balanced by equal and contrary forces, whereas a body is balanced by equal and contrary momenta.

\Article{117}For the same reason, the motion of a solid body differs from the motion of a particle by the mass alone, and thus the equation of the equilibrium or motion of a particle will determine the equilibrium or motion of a solid body, if they be multiplied by its mass.

\Article{118}A moving force is proportional to the quantity of momentum generated by it.

\blank[medium]
\midaligned{\it Reaction equal and contrary to Action.}
\blank[small]

% The following paragraph continues onto original printed page 55.
\Article{119}The law of reaction being equal and contrary to action, is a general induction from observations made on the motions of bodies when placed within certain distances of one another; the law is, that the sum of the momenta generated and estimated in a given direction is zero. It is found by experiment, that if two spheres A and B of the same dimensions and of homogeneous matter, as of gold, be suspended by two threads so as to touch one another when at rest, then if they be drawn aside from the perpendicular to equal heights and let fall at the same instant, they will strike one another centrically, and will destroy each other's motion, so as to remain at rest in the perpendicular. The experiment being repeated with spheres of homogeneous matter, but of different dimensions, if the velocities be inversely as the quantities of matter, the bodies after impinging will remain at rest. It is evident, that in this case, the smaller sphere must descend through a greater space than the larger, in order to acquire the necessary velocity. If the spheres move in the same or in opposite directions, with different momenta, and one strike the other, the body that impinges will lose exactly the quantity of momentum that the other acquires. Thus, in all cases, it is known by experience that reaction is equal and contrary to action, or that equal momenta in opposite directions destroy one another. Daily experience shows that one body cannot acquire motion by the action of another, without depriving the latter body of the same quantity of motion. Iron attracts the magnet with the same force that it is attracted by it; the same thing is seen in electrical attractions and repulsions, and also in animal forces; for whatever may be the moving principle of man and animals, it is found they receive by the reaction of matter, a force equal and contrary to that which they communicate, and in this respect they are subject to the same laws as inanimate beings.

\blank[medium]
\midaligned{\it Mass proportional to Weight.}
\blank[small]

\Article{120}In order to show that the mass of bodies is proportional to their weight, a mode of defining their mass without weighing them must be employed; the experiments that have been described afford the means of doing so, for having arrived at the preceding results, with spheres formed of matter of the same kind, it is found that one of the bodies may be replaced by matter of another kind, but of different dimensions from that replaced. That which produces the same effects as the mass replaced, is considered as containing the same mass or quantity of matter. Thus the mass is defined independent of weight, and as in any one point of the earth's surface every particle of matter tends to move with the same velocity by the action of gravitation, the sum of their tendencies constitutes the weight of a body; hence the mass of a body is proportional to its weight, at one and the same place.

\blank[medium]
\midaligned{\it Density.}
\blank[small]

% The following paragraph continues onto original printed page 56.
\Article{121}Suppose two masses of different kinds of matter, A, of hammered gold, and B of cast copper. If A in motion will destroy the motion of a third mass of matter C, and twice B is required to produce the same effect, then the density of A is said to be double the density of B.

\blank[medium]
\midaligned{\it Mass proportional to the Volume into the Density.}
\blank[small]

\Article{122}The masses of bodies are proportional to their volumes multiplied by their densities; for if the quantity of matter in a given cubical magnitude of a given kind of matter, as water, be arbitrarily assumed as the unit, the quantity of matter in another body of the same magnitude of the density \m{\HeavensRho}, will be represented by \m{\HeavensRho}; and if the magnitude of the second body to that of the first be as \m{m} to 1, the quantity of matter in the second body will be represented by \m{m\times\HeavensRho}.

\blank[medium]
\midaligned{\it Specific Gravity.}
\blank[small]

\Article{123}The densities of bodies of equal volumes are in the ratio of their weights, since the weights are proportional to their masses; therefore, by assuming for the unit of density the maximum density of distilled water at a constant temperature, the density of a body will be the ratio of its weight to that of a like volume of water reduced to this maximum.

\par\noindent\hskip1em\relax This ratio is the specific gravity of a body.

\blank[medium]
\midaligned{\it Equilibrium of two Bodies.}
\blank[small]

\Article{124}If two heavy bodies be attached to the extremities of an inflexible line without mass, which may turn freely on one of its points; when in equilibrio, their masses are reciprocally as their distances from the point of motion.

% This demonstration continues onto original printed page 57.
\testpage[15]
\placefigure[left,none]{}{\useMPgraphic{heavens:figure34}}
\par\noindent\hskip1em\relax {\it Demonstration.}---For, let two heavy bodies \m{m} and \m{m'}, fig.~34, be attached to the extremities of an inflexible line, free to turn round one of its points \m{n}, and suppose the line to be bent in \m{n}, but so little, that \m{m'nm} only differs from two right angles by an indefinitely small angle \m{amn}, which may be represented by \m{\HeavensOmega}. If \m{g} be the force of gravitation, \m{gm}, \m{gm'} will be the gravitation of the two bodies. But the gravitation \m{gm} acting in the direction \m{na} may be resolved into two forces, one in the direction \m{mn}, which is destroyed by the fixed point \m{n}, and another acting on \m{m'} in the direction \m{m'm}. Let \m{mn=f}, \m{m'n=f'}; then \m{m'm=f+f'} very nearly. Hence the whole force \m{gm} is to the part acting on \m{m'\mathrel{::}na:mm'}, and the action of \m{m} on \m{m'}, is \m{\frac{gm(f+f')}{na}}; but \m{m'n:na\mathrel{::}1:\HeavensOmega}, for the arc is so small that it may be taken for its sine. Hence \m{na=\HeavensOmega\mathbin{.}f'}, and the action of \m{m} on \m{m'} is \m{\frac{gm\mathbin{.}(f+f')}{\HeavensOmega f'}}.
\par\flushsidefloats

\par\noindent\hskip1em\relax In the same manner it may be shown that the action of \m{m'} on \m{m} is \m{\frac{gm'(f+f')}{\HeavensOmega f}}; but when the bodies are in equilibrio, these forces must be equal: therefore
\startformula
  \frac{gm(f+f')}{\HeavensOmega f'}=\frac{gm'(f+f')}{\HeavensOmega f},
\stopformula
\noindent whence \m{gm\mathbin{.}f=gm'\mathbin{.}f'}, or \m{gm:gm'\mathrel{::}f':f}, which is the law of equilibrium in the lever, and shows the reciprocal action of parallel forces.

\blank[medium]
\midaligned{\it Equilibrium of a System of Bodies.}
\blank[small]

\Article{125}The equilibrium of a system of bodies may be found, when the system is acted on by any forces whatever, and when the bodies also mutually act on, or attract each other.

\par\noindent\hskip1em\relax {\it Demonstration.}---Let \m{m}, \m{m'}, \m{m''}, \&c., be a system of bodies attracted by a force whose origin is in \m{\mathrm{S}}, fig.~35; and suppose each body to act on all the other bodies, and also to be itself subject to the action of each,---the action of all these forces on the bodies \m{m}, \m{m'}, \m{m''}, \&c., are as the masses of these bodies and the intensities of the forces conjointly.

\placefigure[right,none]{}{\useMPgraphic{heavens:figure35}}
\par\noindent\hskip1em\relax Let the action of the forces on one body, as \m{m}, be first considered; and, for simplicity, suppose the number of bodies to be only three---\m{m}, \m{m'}, and \m{m''}. It is evident that \m{m} is attracted by the force at \m{\mathrm{S}}, and also urged by the reciprocal action of the bodies \m{m'} and \m{m''}.
\par\flushsidefloats

% The following paragraph continues onto original printed page 58.
\par\noindent\hskip1em\relax Suppose \m{m'} and \m{m''} to remain fixed, and that \m{m} is arbitrarily moved to \m{n}: then \m{mn} is the virtual velocity of \m{m}; and if the perpendiculars \m{na}, \m{nb}, \m{nc} be drawn, the lines \m{ma}, \m{mb}, \m{mc}, are the virtual velocities of \m{m} resolved in the direction of the forces which act on \m{m}. Hence, by the principle of virtual velocities, if the action of the force at \m{\mathrm{S}} on \m{m} be multiplied by \m{ma}, the mutual action of \m{m} and \m{m'} by \m{mb}, and the mutual action of \m{m} and \m{m''} by \m{mc}, the sum of these products must be zero when the point \m{m} is in equilibrio; or, \m{m} being the mass, if the action of \m{\mathrm{S}} on \m{m} be \m{\Force\mathbin{.}m}, and the reciprocal actions of \m{m} on \m{m'} and \m{m''} be \m{p}, \m{p'}, then
\startformula
  m\Force\times ma+p\times mb+p'\times mc=0.
\stopformula
\noindent\hskip1em\relax Now, if \m{m} and \m{m''} remain fixed, and that \m{m'} is moved to \m{n'}, then
\startformula
  m'\Force'\times m'a'+p\times m'b'+p''\times m'c'=0.
\stopformula
\noindent And a similar equation may be found for each body in the system. Hence the sum of all these equations must be zero when the system is in equilibrio. If, then, the distances \m{\mathrm{S}m}, \m{\mathrm{S}m'}, \m{\mathrm{S}m''}, be represented by \m{s}, \m{s'}, \m{s''}, and the distances \m{mm'}, \m{mm''}, \m{m'm''}, by \m{f}, \m{f'}, \m{f''}, we shall have
% The source repeats p in the second mutual-action sum and prints a plus/minus
% before "&c."; these are retained rather than silently regularized.
\startformula
  \HeavensSigma\mathpunct{.}m\Force\HeavensDelta s
  +\HeavensSigma\mathpunct{.}p\HeavensDelta f
  +\HeavensSigma\mathpunct{.}p\HeavensDelta f'
  \pm,\ \text{\&c.}=0,
\stopformula
\noindent \m{\HeavensSigma} being the sum of finite quantities; for it is evident that
\startformula
  \HeavensDelta f=mb+m'b',\quad
  \HeavensDelta f'=mc+m''c'',\quad\text{and so on.}
\stopformula
\noindent\hskip1em\relax If the bodies move on surfaces, it is only necessary to add the terms \m{\mathrm{R}\HeavensDelta r}, \m{\mathrm{R}'\HeavensDelta r'}, \&c., in which \m{\mathrm{R}} and \m{\mathrm{R}'} are the pressures or resistances of the surfaces, and \m{\HeavensDelta r}, \m{\HeavensDelta r'} the elements of their directions or the variations of the normals. Hence in equilibrio
\startformula
  \HeavensSigma\mathpunct{.}m\Force\HeavensDelta s
  +\HeavensSigma\mathpunct{.}p\HeavensDelta f
  +\text{\&c.}+\mathrm{R}\HeavensDelta r+\mathrm{R}'\HeavensDelta r',\ \text{\&c.}=0.
\stopformula
\noindent\hskip1em\relax Now, the variation of the normal is zero; consequently the pressures vanish from this equation: and if the bodies be united at fixed distances from each other, the lines \m{mm'}, \m{m'm''}, \&c., or \m{f}, \m{f'}, \&c., are constant:---consequently \m{\HeavensDelta f=0}, \m{\HeavensDelta f'=0}, \&c.

\par\noindent\hskip1em\relax The distance \m{f} of two points \m{m} and \m{m'} in space is
\startformula
  f=\sqrt{(x'-x)^2+(y'-y)^2+(z'-z)^2},
\stopformula
\noindent \m{x}, \m{y}, \m{z}, being the co-ordinates of \m{m}, and \m{x'}, \m{y'}, \m{z'}, those of \m{m'}; so that the variations may be expressed in terms of these quantities: and if they be taken such that \m{\HeavensDelta f=0}, \m{\HeavensDelta f'=0}, \&c., the mutual action of the bodies will also vanish from the equation, which is reduced to
\BookFormula{14}{\HeavensSigma\mathpunct{.}m\Force\mathpunct{.}\HeavensDelta s=0.}

% Article 126 continues onto original printed page 59.
\Article{126}Thus in every case the sum of the products of the forces into the elementary variations of their directions is zero when the system is in equilibrio, provided the conditions of the connexion of the system be observed in their variations or virtual velocities, which are the only indications of the mutual dependence of the different parts of the system on each other.
% Book I, Chapter III, articles 127--134; printed pp. 59--64, PDF pp. 135--140.
% Figures 36--38 are schematic MetaFun reconstructions.

\startuseMPgraphic{heavens:figure36}
  numeric u; u := .22mm;
  pair N,M,Mp,A,Ap;
  N := (0,0); M := (67u,0); Mp := (-100u,0);
  A := (64u,-20u); Ap := (-96u,30u);
  pickup pencircle scaled .35pt;
  draw Mp--M; draw Ap--A;
  draw M{down}..A; draw Mp{up}..Ap;
  label.bot(textext("\tfx N"),N);
  label.rt(textext("\tfx\m{m}"),M);
  label.llft(textext("\tfx\m{m'}"),Mp);
  label.bot(textext("\tfx\m{a}"),A);
  label.top(textext("\tfx\m{a'}"),Ap);
  label(textext("\tfx{\it fig.}\,36."),(18u,29u));
\stopuseMPgraphic

\startuseMPgraphic{heavens:figure37}
  numeric u; u := .21mm;
  pair O,M,N,m,n,A,B,C,D,P;
  O := (0,0); M := (65u,67u); N := (44u,81u);
  m := (65u,-40u); n := (44u,-50u);
  A := (-22u,-40u); B := (-27.5u,-50u);
  C := (44u,0); D := (65u,0); P := (27.5u,-57.8u);
  pickup pencircle scaled .35pt;
  draw (-3u,0)--(114u,0); draw O--(-52u,-95u); draw O--(0,105u);
  draw O--M--N--cycle; draw M--m; draw N--n;
  draw O--m--n--cycle; draw P--m; draw O--P;
  draw m--D; draw n--C;
  draw A--m dashed evenly; draw B--n dashed evenly;
  label.lft(textext("\tfx\m{o}"),O);
  label.rt(textext("\tfx\m{x}"),(114u,0));
  label.bot(textext("\tfx\m{y}"),(-52u,-95u));
  label.lft(textext("\tfx\m{z}"),(0,105u));
  label.rt(textext("\tfx M"),M); label.ulft(textext("\tfx N"),N);
  label.rt(textext("\tfx\m{m}"),m); label.bot(textext("\tfx\m{n}"),n);
  label.lft(textext("\tfx A"),A); label.lft(textext("\tfx B"),B);
  label.top(textext("\tfx C"),C); label.top(textext("\tfx D"),D);
  label.bot(textext("\tfx P"),P);
  label.ulft(textext("\tfx\m{p}"),(58u,0));
  label(textext("\tfx{\it fig.}\,37."),(94u,92u));
\stopuseMPgraphic

\startuseMPgraphic{heavens:figure38}
  numeric u; u := .23mm;
  pair O,A,B,C,D,E,m,n;
  O := (0,70u); A := (0,10u); B := (0,-10u);
  C := (60u,70u); D := (80u,70u); E := (80u,-10u);
  m := (80u,10u); n := (60u,-10u);
  pickup pencircle scaled .35pt;
  draw O--(120u,70u); draw O--(0,-50u);
  draw O--D--E--B; draw A--m; draw C--n;
  draw O--m--n--cycle; draw O--E;
  label.ulft(textext("\tfx\m{o}"),O);
  label.rt(textext("\tfx\m{x}"),(120u,70u));
  label.bot(textext("\tfx\m{y}"),(0,-50u));
  label.lft(textext("\tfx A"),A); label.lft(textext("\tfx B"),B);
  label.top(textext("\tfx C"),C); label.top(textext("\tfx D"),D);
  label.rt(textext("\tfx E"),E);
  label.rt(textext("\tfx\m{m}"),m); label.bot(textext("\tfx\m{n}"),n);
  label(textext("\tfx{\it fig.}\,38."),(85u,94u));
\stopuseMPgraphic

% Original printed page 59; the preceding article ends on this page.
\Article{127}The converse of this law is also true---that when the principle of virtual velocities exists, the system is held in equilibrio by the forces at \m{\mathrm{S}} alone.

\par\noindent\hskip1em\relax {\it Demonstration.}---For if it be not, each of the bodies would acquire a velocity \m{v}, \m{v'}, \&c., in consequence of the forces \m{m\Force}, \m{m'\Force'}, \&c. If \m{\HeavensDelta n}, \m{\HeavensDelta n'}, \&c., be the elements of their direction, then
\startformula
  \HeavensSigma\mathpunct{.}\,m\Force\HeavensDelta s
  -\HeavensSigma\mathpunct{.}\,mv\HeavensDelta n=0.
\stopformula
\noindent The virtual velocities \m{\HeavensDelta n}, \m{\HeavensDelta n'}, \&c., being arbitrary, may be assumed equal to \m{vdt}, \m{v'dt}, \&c., the elements of the space moved over by the bodies; or to \m{v}, \m{v'}, \&c., if the element of the time be unity. Hence
\startformula
  \HeavensSigma\mathpunct{.}\,m\Force\HeavensDelta s
  -\HeavensSigma\mathpunct{.}\,mv^2=0.
\stopformula
\noindent It has been shown that in all cases \m{\HeavensSigma\mathpunct{.}\,m\Force\HeavensDelta s=0}, if the virtual velocities be subject to the conditions of the system. Hence, also, \m{\HeavensSigma\mathpunct{.}\,mv^2=0}; but as all squares are positive, the sum of these squares can only be zero if \m{v=0}, \m{v'=0}, \&c. Therefore the system must remain at rest, in consequence of the forces \m{\Force m}, \&c., alone.

\blank[medium]
\midaligned{\it Rotatory Pressure.}
\blank[medium]

% The original reads "and billiard-ball", without a second indefinite article.
\Article{128}Rotation is the motion of a body, or system of bodies, about a line or point. Thus the earth revolves about its axis, and billiard-ball about its centre.

\Article{129}A rotatory pressure or moment is a force that causes a system of bodies, or a solid body, to rotate about any point or line. It is expressed by the intensity of the motive force or momentum, multiplied by the distance of its direction from the point or line about which the system or solid body rotates.

\blank[medium]
\midaligned{\it On the Lever.}
\blank[medium]

\Article{130}The lever first gave the idea of rotatory pressure or moments, for it revolves about the point of support or fulcrum.

% Original printed page 60 begins within this paragraph.
\par\noindent\hskip1em\relax When the lever \m{mm'}, fig.~36, is in equilibrio, in consequence of forces applied to two heavy bodies at its extremities, the rotatory pressure of these forces, with regard to \m{\mathrm{N}}, the point of support, must be equal and contrary.

\placefigure[left,none]{}{\useMPgraphic{heavens:figure36}}
\par\noindent\hskip1em\relax {\it Demonstration.}---Let \m{ma}, \m{m'a'}, fig.~36, which are proportional to the velocities, represent the forces acting on \m{m} and \m{m'} during the indefinitely small time in which the bodies \m{m} and \m{m'} describe the indefinitely small spaces \m{ma}, \m{m'a'}. The distance of the direction of the forces \m{ma}, \m{m'a'} from the fixed point \m{\mathrm{N}}, are \m{\mathrm{N}m}, \m{\mathrm{N}m'}; and the momentum of \m{m} into \m{\mathrm{N}m}, must be equal to the momentum of \m{m'} into \m{\mathrm{N}m'}; that is, the product of \m{ma} by \m{\mathrm{N}m} and the mass \m{m}, must be equal to the product of \m{m'a'} by \m{\mathrm{N}m'} and the mass \m{m'} when the lever is in equilibrio;
\par\flushsidefloats
\noindent or,
\startformula
  ma\times\mathrm{N}m\times m=m'a'\times\mathrm{N}m'\times m'.
\stopformula
\noindent But \m{ma\times\mathrm{N}m} is twice the triangle \m{\mathrm{N}ma}, and \m{m'a'\times\mathrm{N}m'} is twice the triangle \m{\mathrm{N}m'a'}; hence twice the triangle \m{\mathrm{N}ma} into the mass \m{m}, is equal to twice the triangle \m{\mathrm{N}m'a'} into the mass \m{m'}, and these are the rotatory pressures which cause the lever to rotate about the fulcrum; thus, in equilibrio, the rotatory pressures are equal and contrary, and the moments are inversely as the distances from the point of support.

\blank[medium]
\midaligned{\it Projection of Lines and Surfaces.}
\blank[medium]

\placefigure[left,none]{}{\useMPgraphic{heavens:figure37}}
\Article{131}Surfaces and areas may be projected on the co-ordinate planes by letting fall perpendiculars from every point of them on these planes. For let \m{o\mathrm{MN}}, fig.~37, be a surface meeting the plane \m{xoy} in \m{o}, the origin of the co-ordinates, but rising above it towards \m{\mathrm{MN}}. If perpendiculars be drawn from every point of the area \m{o\mathrm{MN}} on the plane \m{xoy}, they will trace the area \m{omn}, which is the projection of \m{o\mathrm{MN}}.

\par\noindent\hskip1em\relax Since, by hypothesis, \m{xoy} is a right angle, if the lines \m{m\mathrm{D}}, \m{n\mathrm{C}}, be drawn parallel to \m{oy}, \m{\mathrm{DC}} is the projection of \m{mn} on the axis \m{ox}. In the same manner \m{\mathrm{AB}} is the projection of the same line on \m{oy}.
\par\flushsidefloats

% Original printed page 61.
\blank[medium]
\midaligned{\it Equilibrium of a System of Bodies invariably united.}
\blank[medium]

\Article{132}A system of bodies invariably united will be in equilibrio upon a point, if the sum of the moments of rotation of all the forces that act upon it vanish, when estimated parallel to three rectangular co-ordinates.

\par\noindent\hskip1em\relax {\it Demonstration.}---Suppose a system of bodies invariably united, moving about a fixed point \m{o} in consequence of an impulse and a force of attraction; \m{o} being the origin of the attractive force and of the co-ordinates.

\par\noindent\hskip1em\relax Let one body be considered at a time, and suppose it to describe the indefinitely small arc \m{\mathrm{MN}}, fig.~37, in an indefinitely small time, and let \m{mn} be the projection of this arc on the plane \m{xoy}. If \m{m} be the mass of the body, then \m{m\times mn} is its momentum, estimated in the plane \m{xoy}; and if \m{o\mathrm{P}} be perpendicular to \m{mn}, it is evident that \m{m\times mn\times o\mathrm{P}} is its rotatory pressure. But \m{mn\times o\mathrm{P}} is twice the triangle \m{mon}; hence the rotatory pressure is equal to the mass \m{m} into twice the triangle \m{mon} that the body could describe in an element of time. But when \m{m} is at rest, the rotatory pressure must be zero; hence in equilibrio, \m{m\times mn\times o\mathrm{P}=0}.

% Source reads "fig. 33", although the adjacent diagram is numbered 38.
\placefigure[right,none]{}{\useMPgraphic{heavens:figure38}}
\par\noindent\hskip1em\relax Let \m{omn}, fig.~33, be the projected area, and complete the parallelogram \m{o\mathrm{DEB}}; then if \m{o\mathrm{D}}, \m{o\mathrm{A}}, the co-ordinates of \m{m}, be represented by \m{x} and \m{y}, it is evident that \m{y} increases, while \m{x} diminishes; hence
\par\flushsidefloats
\startformula
  \mathrm{CD}=-dx,\quad\text{and}\quad\mathrm{AB}=dy.
\stopformula
% The source expresses these areas as half a side alone; this notation is retained.
\noindent Join \m{o\mathrm{E}}, then
\startformula
  no\mathrm{E}=\tfrac12 n\mathrm{D},
\stopformula
\noindent because the triangle and parallelogram are on the same base and between the same parallels; also \m{mo\mathrm{E}=\tfrac12\mathrm{AE}}: hence the triangle
\startformula
  mon=\tfrac12\{n\mathrm{D}+\mathrm{AE}.\}
\stopformula
\noindent Now \m{n\mathrm{D}=-dx(y+dy)} and \m{\mathrm{AE}=xdy}, therefore
\startformula
  mon=\tfrac12(xdy-ydx)-\tfrac12 dxdy;
\stopformula
\noindent but when the arc \m{mn} is indefinitely small, \m{\tfrac12 dxdy=\tfrac12 n\mathrm{E}\mathbin{.}m\mathrm{E}} may be omitted in comparison of the first powers of these quantities, hence the triangle
\startformula
  mon=\tfrac12(xdy-ydx),
\stopformula
% Original printed page 62 begins within this paragraph.
\noindent therefore \m{m(xdy-ydx)=0} is the rotatory pressure in the plane \m{xoy} when \m{m} is in equilibrio. A similar equation must exist for each co-ordinate plane when \m{m} is in a state of equilibrium with regard to each axis, therefore also
\startformula
  m(xdz-zdx)=0,\quad m(ydz-zdy)=0.
\stopformula
\noindent The same may be proved for every body in the system, consequently when the whole is in equilibrio on the point \m{o}
\BookFormula{15}{
  \startalign[n=1,align=middle]
    \NC \HeavensSigma m(xdy-ydx)=0 \qquad \HeavensSigma m(xdz-zdx)=0 \NR
    \NC \HeavensSigma m(ydz-zdy)=0. \NR
  \stopalign
}

\Article{133}This property may be expressed by means of virtual velocities, namely, that a system of bodies will be at rest, if the sum of the products of their momenta by the elements of their directions be zero, or by article 125
\startformula
  \HeavensSigma m\Force\HeavensDelta s=0.
\stopformula
\noindent Since the mutual distances of the parts of the system are invariable, if the whole system be supposed to be turned by an indefinitely small angle about the axis \m{oz}, all the co-ordinates \m{z'}, \m{z''}, \&c., will be invariable. If \m{\HeavensDelta\HeavensVarpi} be any arbitrary variation, and if
\startformula
  \startalign[n=1,align=middle]
    \NC \HeavensDelta x=y\HeavensDelta\HeavensVarpi \qquad
        \HeavensDelta y=-x\HeavensDelta\HeavensVarpi \NR
    \NC \HeavensDelta x'=y'\HeavensDelta\HeavensVarpi \qquad
        \HeavensDelta y'=-x'\HeavensDelta\HeavensVarpi; \NR
  \stopalign
\stopformula
\noindent then \m{f} being the mutual distance of the bodies \m{m} and \m{m'} whose co-ordinates are \m{x,y,z}; \m{x',y',z'}, there will arise
\startformula
  \HeavensDelta f=\HeavensDelta\sqrt{(x'-x)^2+(y'-y)^2+(z'-z)^2}=
\stopformula
\startformula
  \frac{x'-x}{f}(\HeavensDelta x'-\HeavensDelta x)
  +\frac{y'-y}{f}(\HeavensDelta y'-\HeavensDelta y)=
\stopformula
% The printed prefactor in the final line is 1/2; retained as printed.
\startformula
  \tfrac12\{(x'-x)(y'-y)\HeavensDelta\HeavensVarpi
    -(y'-y)(x'-x)\HeavensDelta\HeavensVarpi\}=0.
\stopformula
\noindent So that the values assumed for \m{\HeavensDelta x}, \m{\HeavensDelta y}, \m{\HeavensDelta x'}, \m{\HeavensDelta y'} are not incompatible with the invariability of the system. It is therefore a permissible assumption.

\par\noindent\hskip1em\relax Now if \m{s} be the direction of the force acting on \m{m}, its variation is
\startformula
  \HeavensDelta s=\Der{s}{x}\HeavensDelta x+\Der{s}{y}\HeavensDelta y,
\stopformula
\noindent since \m{z} is constant; and substituting the preceding values of \m{\HeavensDelta x}, \m{\HeavensDelta y}, the result is
\startformula
  \HeavensDelta s=\Der{s}{x}\mathbin{.}y\HeavensDelta\HeavensVarpi
    -\Der{s}{y}\mathbin{.}x\HeavensDelta\HeavensVarpi
    =\HeavensDelta\HeavensVarpi\left\{\Der{s}{x}\mathbin{.}y-\Der{s}{y}x\right\}
\stopformula
\noindent or, multiplying by the momentum \m{\Force m},
\startformula
  \Force m\HeavensDelta s=\Force m\left\{y\Der{s}{x}-x\Der{s}{y}\right\}\HeavensDelta\HeavensVarpi.
\stopformula

% Original printed page 63.
\noindent In the same manner with regard to the body \m{m'}
\startformula
  \Force'm'\HeavensDelta s'=\Force'm'\left\{y'\Der{s'}{x'}-x'\Der{s'}{y'}\right\}\HeavensDelta\HeavensVarpi,
\stopformula
\noindent and so on; and thus the equation \m{\HeavensSigma m\Force\HeavensDelta s=0} becomes
\startformula
  \HeavensSigma m\Force\left\{y\Der{s}{x}-x\Der{s}{y}\right\}=0.
\stopformula
\noindent It follows, from the same reasoning, that
\startformula
  \HeavensSigma m\Force\left\{z\Der{s}{x}-x\Der{s}{z}\right\}=0,
\stopformula
\startformula
  \HeavensSigma m\Force\left\{z\Der{s}{y}-y\Der{s}{z}\right\}=0.
\stopformula
\noindent In fact, if \m{\ForceX}, \m{\ForceY}, \m{\ForceZ} be the components of the force \m{\Force} in the direction of the three axes, it is evident that
\startformula
  \ForceX=\Force\Der{s}{x};\quad
  \ForceY=\Force\Der{s}{y};\quad
  \ForceZ=\Force\Der{s}{z};
\stopformula
\noindent and these equations become
\BookFormula{16}{
  \startalign[n=1,align=middle]
    \NC \HeavensSigma my\mathbin{.}\ForceX-\HeavensSigma mx\mathbin{.}\ForceY=0 \NR
    \NC \HeavensSigma mz\mathbin{.}\ForceX-\HeavensSigma mx\mathbin{.}\ForceZ=0 \NR
    \NC \HeavensSigma mz\mathbin{.}\ForceY-\HeavensSigma my\mathbin{.}\ForceZ=0 \NR
  \stopalign
}
% Source says "round x" here, although the surrounding discussion concerns the z moment.
\noindent But \m{\HeavensSigma m\Force y\Der{s}{x}} expresses the sum of the moments of the forces parallel to the axis of \m{x} to turn the system round that of \m{z}, and \m{\HeavensSigma m\Force x\Der{s}{y}} that of the forces parallel to the axis of \m{y} to do the same, but estimated in the contrary direction;---and it is evident that the forces parallel to \m{z} have no effect to turn the system round \m{x}. Therefore the equation \m{\HeavensSigma m\Force\left(y\Der{s}{x}-x\Der{s}{y}\right)=0}, expresses that the sum of the moments of rotation of the whole system relative to the axis of \m{z} must vanish, that the equilibrium of the system may subsist. And the same being true for the other rectangular axes (whose positions are arbitrary), there results this general theorem, viz., that in order that a system of bodies may be in equilibro upon a point, the sum of the moments of rotation of all the forces that act on it must vanish when estimated parallel to any three rectangular co-ordinates.

% Original printed page 64 begins within this paragraph.
\Article{134}These equations are sufficient to ensure the equilibrium of the system when \m{o} is a fixed point; but if \m{o}, the point about which it rotates, be not fixed, the system, as well as the origin \m{o}, may be carried forward in space by a motion of translation at the same time that the system rotates about \m{o}, like the earth, which revolves about the sun at same time that it turns on its axis. In this case it is not only necessary for the equilibrium of the system that its rotatory pressure should be zero, but also that the forces which cause the translation when resolved in the direction of the axis \m{ox}, \m{oy}, \m{oz}, should be zero for each axis separately.
% Book I, Chapter III, articles 135--143; printed pages 64--68.
% Original PDF pages 140--144. Figures 39--40 are schematic reconstructions.

\startuseMPgraphic{heavens:figure39}
  numeric u; u := .23mm;
  pair O;
  O := (65u,0);
  pickup pencircle scaled .35pt;
  draw (0,0)--(150u,0);
  draw (80u,-49u)--(80u,58u);
  draw (19u,0)--(19u,30u);
  draw (39u,0)--(39u,48u);
  draw (58u,0)--(58u,-44u);
  draw (93u,0)--(93u,37u);
  draw (111u,0)--(111u,-42u);
  draw (129u,0)--(129u,46u);
  label.lft(textext("\tfx\m{x}"),(0,0));
  label.rt(textext("\tfx\m{x'}"),(150u,0));
  label.top(textext("\tfx\m{z}"),(80u,58u));
  label.bot(textext("\tfx\m{z'}"),(80u,-49u));
  label.top(textext("\tfx\m{o}"),O);
  label.bot(textext("\tfx\m{a}"),(19u,0));
  label.bot(textext("\tfx\m{b}"),(39u,0));
  label.top(textext("\tfx\m{c}"),(58u,0));
  label.bot(textext("\tfx\m{d}"),(93u,0));
  label.top(textext("\tfx\m{e}"),(111u,0));
  label.bot(textext("\tfx\m{f}"),(129u,0));
  label.top(textext("\tfx\m{m}"),(19u,30u));
  label.top(textext("\tfx\m{m'}"),(39u,48u));
  label.bot(textext("\tfx\m{m''}"),(58u,-44u));
  label.top(textext("\tfx\m{m'''}"),(93u,37u));
  label.bot(textext("\tfx\m{m^{\rm iv}}"),(111u,-42u));
  label.top(textext("\tfx\m{m^{\rm v}}"),(129u,46u));
  label(textext("\tfx{\it fig.}\,39."),(126u,78u));
\stopuseMPgraphic

\startuseMPgraphic{heavens:figure40}
  numeric u; u := .24mm;
  pair O,A,M,N;
  O := (0,0); A := (76u,0); M := (76u,76u); N := (76u,-44u);
  pickup pencircle scaled .35pt;
  draw O--(150u,0);
  draw M--N;
  draw (35u,0)--(35u,54u)--(76u,54u);
  draw (53u,0)--(53u,31u)--(76u,31u);
  draw (101u,0)--(101u,20u)--(76u,20u);
  label.lft(textext("\tfx\m{o}"),O);
  label.rt(textext("\tfx\m{x}"),(150u,0));
  label.top(textext("\tfx M"),M);
  label.bot(textext("\tfx N"),N);
  label.llft(textext("\tfx A"),A);
  label.bot(textext("\tfx\m{p}"),(35u,0));
  label.bot(textext("\tfx\m{p'}"),(53u,0));
  label.bot(textext("\tfx\m{p''}"),(101u,0));
  label.urt(textext("\tfx\m{a}"),(76u,54u));
  label.rt(textext("\tfx\m{b}"),(76u,31u));
  label.rt(textext("\tfx\m{c}"),(76u,44u));
  label.ulft(textext("\tfx\m{d}"),(76u,20u));
  label.top(textext("\tfx\m{m}"),(35u,54u));
  label.ulft(textext("\tfx\m{m'}"),(53u,31u));
  label.urt(textext("\tfx\m{m''}"),(101u,20u));
  label(textext("\tfx{\it fig.}\,40."),(130u,54u));
\stopuseMPgraphic

% Original printed page 64, after the conclusion of article 134.
\blank[medium]
\midaligned{\it On the Centre of Gravity.}
\blank[medium]

\Article{135}If the bodies \m{m}, \m{m'}, \m{m''}, \&c., be only acted on by gravity, its effect would be the same on all of them, and its direction may be considered the same also; hence
\startformula
  \Force=\Force'=\Force''=\text{\&c.},
\stopformula
\noindent and also the directions
\startformula
  \Der{s}{x}=\Der{s}{x'}=\text{\&c.}\quad
  \Der{s}{y}=\Der{s}{y'}=\text{\&c.}\quad
  \Der{s}{z}=\Der{s}{z'}=\text{\&c.},
\stopformula
\noindent are the same in this case for all the bodies, so that the equations of rotatory pressure become
\startformula
  \startalign[n=1,align=middle]
  \NC \Force\left\{\Der{s}{x}\mathbin{.}\HeavensSigma my-\Der{s}{y}\HeavensSigma mx\right\}=0 \NR
  \NC \Force\left\{\Der{s}{z}\mathbin{.}\HeavensSigma my-\Der{s}{y}\HeavensSigma mz\right\}=0 \NR
  \NC \Force\left\{\Der{s}{x}\mathbin{.}\HeavensSigma mz-\Der{s}{z}\HeavensSigma mx\right\}=0 \NR
  \stopalign
\stopformula
\noindent or, if \m{\ForceX}, \m{\ForceY}, \m{\ForceZ}, be considered as the components of gravity in the three co-ordinate axes by article 133
\BookFormula{17}{
  \startalign[n=1,align=middle]
  \NC \ForceX\mathbin{.}\HeavensSigma my-\ForceY\mathbin{.}\HeavensSigma mx=0 \NR
  \NC \ForceZ\mathbin{.}\HeavensSigma my-\ForceY\mathbin{.}\HeavensSigma mz=0 \NR
  \NC \ForceX\mathbin{.}\HeavensSigma mz-\ForceZ\mathbin{.}\HeavensSigma mx=0 \NR
  \stopalign
}
\noindent It is evident that these equations will be zero, whatever the direction of gravity may be, if
\BookFormula{18}{
  \HeavensSigma mx=0,\quad \HeavensSigma my=0,\quad \HeavensSigma mz=0.
}
\noindent\hskip1em\relax Now since \m{\Force\Der{s}{x}}, \m{\Force\Der{s}{y}}, \m{\Force\Der{s}{z}}, are the components of the force of gravity in the three co-ordinates \m{ox}, \m{oy}, \m{oz},
\startformula
  \Force\mathbin{.}\Der{s}{x}\mathbin{.}\HeavensSigma m;\quad
  \Force\mathbin{.}\Der{s}{y}\mathbin{.}\HeavensSigma m;\quad
  \Force\mathbin{.}\Der{s}{z}\mathbin{.}\HeavensSigma m;
\stopformula
\noindent are the forces which translate the system parallel to these axes. But if \m{o} be a fixed point, its reaction would destroy these forces.

% Original printed page 65 begins in the preceding paragraph.
\noindent By article 49,
\startformula
  \left(\Der{s}{x}\right)^2+\left(\Der{s}{y}\right)^2+\left(\Der{s}{z}\right)^2=1
\stopformula
\noindent is the diagonal of a parallelopiped, of which
\startformula
  \Der{s}{x},\quad\Der{s}{y},\quad\Der{s}{z},
\stopformula
\noindent are the sides; therefore these three compose one resulting force equal to \m{\Force\mathbin{.}\HeavensSigma m}. This resulting force is the weight of the system which is thus resisted or supported by the reaction of the fixed point \m{o}.

\Article{136}The point \m{o} round which the system is in equilibrio, is the centre of gravity of the system, and if that point be supported, the whole will be in equilibrio.

\blank[big]
\midaligned{\it On the Position and Properties of the Centre of Gravity.}
\blank[medium]

\Article{137}It appears from the equations (18), that if any plane passes through the centre of gravity of a system of bodies, the sum of the products of the mass of each body by its distance from that plane is zero. For, since the axes of the co-ordinates are arbitrary, any one of them, as \m{xox'}, fig.~39, may be assumed to be the section of the plane in question, the centre of gravity of the system of bodies \m{m}, \m{m'}, \&c., being in \m{o}.
\placefigure[right,none]{}{\useMPgraphic{heavens:figure39}}
\noindent If the perpendiculars \m{ma}, \m{m'b}, \&c., be drawn from each body on the plane \m{xox'}, the product of the mass \m{m} by the distance \m{ma} plus the product of \m{m'} by \m{m'b} plus, \&c., must be zero; or, representing the distances by \m{z}, \m{z'}, \m{z''}, \&c., then
\par\flushsidefloats
% The source omits the prime on z in the second term; retained verbatim.
\startformula
  mz+m'z-m''z''+m'''z'''+\text{\&c.}=0;
\stopformula
\noindent or, according to the usual notation,
\startformula
  \HeavensSigma\mathbin{.}mz=0.
\stopformula
\noindent And the same property exists for the other two co-ordinate planes. Since the position of the co-ordinate planes is arbitrary, the property obtains for every set of co-ordinate planes having their origin in \m{o}. It is clear that if the distances \m{ma}, \m{m'b}, \&c., be positive on one side of the plane, those on the other side must be negative, otherwise the sum of the products could not be zero.

% Original printed page 66 begins within the preceding paragraph.
\Article{138}When the centre of gravity is not in the origin of the co-ordinates, it may be found if the distances of the bodies \m{m}, \m{m'}, \m{m''}, \&c., from the origin and from each other be known.
\par\blank[small]
\midaligned{\useMPgraphic{heavens:figure40}}
\blank[small]
\noindent\hskip1em\relax {\it Demonstration.}---For let \m{o}, fig.~40, be the origin, and \m{c} the centre of gravity of the system \m{m}, \m{m'}, \&c. Let \m{\mathrm{MN}} be the section of a plane passing through \m{c}; then by the property of the centre of gravity just explained,
\par\flushsidefloats
\startformula
  m\mathbin{.}ma+m'\mathbin{.}m'b-m''\mathbin{.}m''d+\text{\&c.}=0;
\stopformula
\noindent but
\startformula
  ma=o\mathrm{A}-op;\quad m'b=o\mathrm{A}-op',\quad\text{\&c.\ \&c.},
\stopformula
\noindent hence
\startformula
  m(o\mathrm{A}-op)+m'(o\mathrm{A}-op')+\text{\&c.}=0;
\stopformula
\noindent or if \m{\mathrm{A}o} be represented by \m{\bar{x}}, and \m{op}, \m{op'}, \m{op''}, \&c., by \m{x}, \m{x'}, \m{x''}, \&c., then will
% The minus signs on the third mass are retained from the original, including
% the inconsistent signs in equation (19), rather than silently corrected.
\startformula
  m(\bar{x}-x)+m'(\bar{x}-x')-m''(\bar{x}-x'')+\text{\&c.}=0.
\stopformula
\noindent Whence
\startformula
  \bar{x}(m+m'-m''+\text{\&c.})=mx+m'x'-m''x''+\text{\&c.},
\stopformula
\BookFormula{19}{
  \text{and, }\bar{x}=\frac{mx+m'x'+\text{\&c.}}{m+m'-m''+\text{\&c.}}
    =\frac{\HeavensSigma\mathbin{.}mx}{\HeavensSigma m}.
}
\noindent\hskip1em\relax Thus, if the masses of the bodies and their respective distances from the origin of the co-ordinates be known, this equation will give the distance of the centre of gravity from the plane \m{yoz}. In the same manner its distances from the other two co-ordinate planes are found to be
\BookFormula{20}{
  \bar{y}=\frac{\HeavensSigma\mathbin{.}my}{\HeavensSigma m}\qquad
  \bar{z}=\frac{\HeavensSigma\mathbin{.}mz}{\HeavensSigma m}.
}

\Article{139}Thus, because the centre of gravity is determined by its three co-ordinates \m{\bar{x}}, \m{\bar{y}}, \m{\bar{z}}, it is a single point.

% Original printed page 67.
\Article{140}But these three equations give
\startformula
  \bar{x}^2+\bar{y}^2+\bar{z}^2
    =\frac{(\HeavensSigma mx)^2+(\HeavensSigma my)^2+(\HeavensSigma mz)^2}{(\HeavensSigma m)^2},
\stopformula
\noindent or
\startformula
  \startalign[n=1,align=middle]
  \NC \bar{x}^2+\bar{y}^2+\bar{z}^2
    =\frac{\HeavensSigma m(x^2+y^2+z^2)}{\HeavensSigma m} \NR
  \NC {}-\frac{\HeavensSigma mm'\{(x'-x)^2+(y'-y)^2+(z'-z)^2\}}{(\HeavensSigma m)^2} \NR
  \stopalign
\stopformula
\noindent The last term of the second member is the sum of all the products similar to those under \m{\HeavensSigma} when all the bodies of the system are taken in pairs.

\Article{141}It is easy to show that the two preceding values of \m{\bar{x}^2+\bar{y}^2+\bar{z}^2} are identical, or that
\startformula
  \frac{(\HeavensSigma mx)^2}{(\HeavensSigma m)^2}
    =\frac{\HeavensSigma mx^2}{\HeavensSigma m}
      -\frac{\HeavensSigma mm'(x'-x)^2}{(\HeavensSigma m)^2}
\stopformula
\noindent or
\startformula
  (\HeavensSigma mx)^2=\HeavensSigma m\mathbin{.}\HeavensSigma mx^2-\HeavensSigma mm'(x'-x)^2.
\stopformula
\noindent Were there are only two planets, then
% The source reads "Were there are only two planets", retained verbatim.
\startformula
  \HeavensSigma m=m+m',\quad
  \HeavensSigma mx=mx+m'x',\quad
  \HeavensSigma mm'=mm';
\stopformula
\noindent consequently
\startformula
  (\HeavensSigma mx)^2=(mx+m'x')^2=m^2x^2+m'^2x'^2+2mm'xx'.
\stopformula
\noindent With regard to the second member
\startformula
  \startalign[n=1,align=middle]
  \NC \HeavensSigma m\mathbin{.}\HeavensSigma mx^2=(m+m')(mx^2+m'x'^2) \NR
  \NC {}=m^2x^2+m'^2x'^2+mm'x^2+mm'x'^2, \NR
  \stopalign
\stopformula
\noindent and
\startformula
  \HeavensSigma mm'(x'-x)^2=mm'x'^2+mm'x^2-2mm'xx';
\stopformula
\noindent consequently
\startformula
  \startalign[n=1,align=middle]
  \NC \HeavensSigma m\mathbin{.}\HeavensSigma mx^2-\HeavensSigma mm'(x'-x)^2 \NR
  \NC {}=m^2x^2+m'^2x'^2+2mm'xx'=(\HeavensSigma mx)^2. \NR
  \stopalign
\stopformula
\noindent This will be the case whatever the number of planets may be; and as the equations in question are symmetrical with regard to \m{x}, \m{y}, and \m{z}, their second members are identical.

\par\noindent\hskip1em\relax Thus the distance of the centre of gravity from a given point may be found by means of the distances of the different points of the system from this point, and of their mutual distances.

\Article{142}By estimating the distance of the centre of gravity from any three fixed points, its position in space will be determined.

\blank[big]
\midaligned{\it Equilibrium of a Solid Body.}
\blank[medium]

\Article{143}If the bodies \m{m}, \m{m'}, \m{m''}, \&c., be indefinitely small, infinite in number, and permanently united together, they will form a solid mass, whose equilibrium may be determined by the preceding equations.

% Original printed page 68.
\par\noindent\hskip1em\relax For if \m{x}, \m{y}, \m{z}, be the co-ordinates of any one of its indefinitely small particles \m{dm}, and \m{\ForceX}, \m{\ForceY}, \m{\ForceZ}, the forces urging it in the direction of these axes, the equations of its equilibrium will be
\startformula
  \int\ForceX dm=0\qquad \int\ForceY dm=0\qquad \int\ForceZ dm=0
\stopformula
\startformula
  \startalign[n=1,align=middle]
  \NC \int(\ForceX y-\ForceY x)\,dm=0;\quad
      \int(\ForceX z-\ForceZ x)\,dm=0; \NR
  \NC \int(\ForceZ y-\ForceY z)\,dm=0. \NR
  \stopalign
\stopformula
\noindent The three first are the equations of translation, which are destroyed when the centre of gravity is a fixed point; and the last three are the sums of the rotatory pressures.

\par\flushsidefloats
\blank[small]
\startalignment[middle]
  \blackrule[width=45mm,height=.3pt]
\stopalignment
\ifdefined\MechanismCombined\else\expandafter\stoptext\fi
